\textbf{Problème 4.9 (Dujella, 2026).} Soit $k \geqslant 3$ un entier. Est-il vrai que tous les points entiers de la courbe elliptique $y^2=((k-1)x+1)((k+1)x+1)(4kx+1)$ sont donnés par $(x,y) \in \{(0,\pm 1),(16k^3-4k,\pm(128k^6-112k^4+20k^2-1))\}$ ?

\medskip

\textit{Remarque.} La réponse positive est connue si le rang de la courbe vaut 1 (Dujella, Acta Arith. \textbf{94} (2000), 87--101) ou si $k \leqslant 5000$ (Najman, Rocky Mountain J. Math. \textbf{40} (2010), 1979--2002). Également si $k=3n^2+2n-2$ ou $k=\frac{1}{2}(3m^2+5m)$ et que le rang vaut 2 (Dujella, Acta Arith. \textbf{94} (2000), 87--101).
